(1) 作用量为\(\mathcal S=\int dt\int dx[\mu y_t^2/2-Ty_x^2/2-EIy_{xx}^2/2]\)。对\(y\)变分并分部积分两次，体项给\(-\mu y_{tt}+Ty_{xx}-EIy_{xxxx}=0\)，故运动方程\(\mu y_{tt}+EIy_{xxxx}-Ty_{xx}=0\)。拉力项抵抗横向曲率，弯曲项对应四阶空间导数。
(2) 代入\(y=Ae^{i(kx-\omega t)}\)，有\(\mu\omega^2=EIk^4+Tk^2\)。正向波相速度\(v_p=\omega/k=\sqrt{(EI k^2+T)/\mu}\)；群速度\(v_g=d\omega/dk=(2EIk^3+Tk)/(\mu\omega)\)。当弯曲项主导时\(v_g\simeq2v_p\)，不能把此色散波的能流速度简单设为其相速度。
(3) 两种势能量级在\(EIk^4\sim Tk^2\)时平衡，故\(k_*=\sqrt{T/(EI)}\)（\(T=0\)时无交叉尺度）。若\(k\ll k_*\)且\(T>0\)，\(\omega\simeq\sqrt{T/\mu}|k|\)，接近非色散张紧弦波；若\(k\gg k_*\)，\(\omega\simeq\sqrt{EI/\mu}\,k^2\)，弯曲色散成为主导。
(4) 变分后的边界项为\([(-Ty_x+EI y_{xxx})\delta y-EIy_{xx}\delta y_x]\)的符号约定等价形式。固支端\(x=0\)几何条件\(y(0,t)=y_x(0,t)=0\)；自由端\(x=L\)独立允许\(\delta y,\delta y_x\)变化，必须满足\(EIy_{xx}(L,t)=0\)与\(EIy_{xxx}(L,t)-Ty_x(L,t)=0\)，分别对应端弯矩为零和有效端剪力为零。任何仅给\(y_{xx}=y_{xxx}=0\)的自由端条件都只适用于\(T=0\)。
